\documentclass{article}
\usepackage{amssymb, amsmath}
\usepackage[left=1.5cm, right=2cm]{geometry}
\begin{document}

\textbf{Fixing Mutual Authentication}

Encryption does not ensure authenticity: use MAC to authenticate messages\\
Prevent redirection: identify party in challenge\\
Better to use separate keys for each direction\\
Prevent replay and reorder\begin{itemize}
	\item Identify flow and connection
	\item Prevent use of old challenge: randomness, time, or state
\end{itemize}
Do not provide the adversary with oracle access! (like in SNA)\begin{itemize}
	\item Do not compute values from Adversary
	\item Include self-chosen nonce in the protected reply
\end{itemize}

\textbf{Secure Two-Party Handshake Protocol (2PP)}

Use MAC rather than encryption to auth\\
Prevent redirection: include identities $(A,B)$\\
Prevent replay and reorder:\begin{itemize}
	\item Nonces $(N_A,N_B)$
	\item Separate $2^{nd}$ and $3^{rd}$ flows: 3 vs 2 input blocks
\end{itemize}
Provably secure, since MACs are deterministic\\
New nonces are generated for each flow, prevents concurrent attack\\
Flow numbers ensure prevention of redirection

\textbf{Authenticated Request-Control Protocols}

Besides authenticating entities, these protocols auth the exchange of a request and a response between the entities.\\
Required properties:\begin{itemize}
	\item Request Authentication: the request was indeed sent by the sender.
	\item Response Auth: the response was indeed sent by the receiver (to which the request was intended for)
	\item No Replay
\end{itemize}
Five variants:\begin{itemize}
	\item 2PP-RR (two party protocol, request-response)
	\item 2RT-2PP (2 round trip) (a round trip is two flows)
	\item Counter-based-RR
	\item Time-based RR
	\item Key-exchange (session key)
\end{itemize}

\textbf{2PP-RR}

A three-flow nonce-based protocol.\\
Significant drawback:\begin{itemize}
	\item The request is sent by the responder and the initiator sends the response.
	\item So, the initiator must wait for a request rather than sending it
\end{itemize}

\textbf{2RT-PP}

A four-flow nonce-based protocol.\\
Mainly fixes the drawback of 2PP-RR

\textbf{Counter-Based Authenticated RR}

Simple stateful (counter) solution, requiring only one round:\begin{itemize}
	\item Unidrectional (if bidirectional is needed, a separate instance of the protocol for each direction needs to be executed)
	\item Parties maintain synchronized counter $i$ of requests (and responses) to avoid replay attacks
	\item Recipient validates if counter received is $i+1$
	\item Both parties must remember the counter
\end{itemize}
Counters are only initialized once

\textbf{Time-Based Auth RR}

Simple stateful (time) solution, requiring only one round:\begin{itemize}
	\item Use local clocks $T_A,T_B$ instead of counters with two assumptions: bounded delays and bounded clock skews
	\item Responder (Bob):\begin{itemize}
			\item Rejects request if $T_B>T_A+\Delta$ OR if he received a larger $T_A$ already (where $\Delta=\Delta_{skew}+\Delta_{delay}$)
			\item Maintains last $T_A$ received until $T_A+\Delta$
		\end{itemize}
	\item Initiator does not need any state, when can receiver discard this?
\end{itemize}
Don't reuse keys across parties (TWO party protocol)

\textbf{2RT-2PP with Confidentiality}

Secure connection: authentication, freshness, and secrecy\begin{itemize}
	\item Independent keys: for encryption $k.e$, for auth $k.a$
	\item How can we derive them both from a single key $k$? The PRF from before: $k.e=PRF_k("Encrypt"),\ k.a=PRF_k("MAC")$
\end{itemize}
Improve security by changing keys between sessions?

\textbf{2PP Key Exchange Protocol}

Allows generating independent session keys (i.e. each session will have its own key denoted as $k_i^S$\\
These keys are derived from the shared key that we now call a long-term 'master key' denoted as $k^M$\\
Improves security:\begin{itemize}
	\item Exposure of a session key does not expose the master key, or keys of other sessions
	\item Thus, limited amount of ciphertext is exposed if a session key is exposed (only that session, rather than all the ciphertexts of other sessions)
\end{itemize}
$A,N_{A.i}$\\
$N_{A.i},N_{B.i}, PRF_{k^M}(2||'A\leftarrow B||N_{A.i}||N_{B.i})$\\
$N_{B.i},PRF_{k^M}(3||'A\rightarrow B'||N_{A.i}||N_{B.i})$\\
Session key: $k_i^S = PRF_{k^M}(N_{A.i}||N_{B.i})$\\
Use session key to generate MAC and Encryption key\\
A secure PRF is also a secure MAC, this example uses only one cryptographic primitive (the PRF, for session key and handshake).\\
You can't use the same key against two different crypto primitives

\textbf{Key Distribution Center (KDCs)}: Establish a shared key b/w two or more entities, usually with the help of a trusted 3rd party reference

\textbf{KDC}\\
Has huge trust assumption based on centralization\\
Will focus on three party protocols: Alice, Bob, and KDC\\
KDC: shares long-term master keys with all parties (each party will have a key for encryption and another for MAC)\begin{itemize}
	\item Denoted as $k_A^M,k_B^M,...$ for MAC and $k_A^E,k_B^E,...$ for encryption
\end{itemize}
Goal: help parties (A, B) to \underline{establish} a shared master key $k_{AB}$, based on which the parties generate two keys for MAC and encryption, namely $k^M_{AB}$ and $k^E_{AB}$\\
We will study simplified versions of Kerberos (secure) and GSM (insecure)

\textbf{The Kerberos KDC Protocol}\\
KDC Shares keys $k_A^E$ (enc), $k_A^M$ (MAC) with Alice, and $k^E_B,k^M_B$ with Bob\\
Goal: Alice and Bob share $k_{AB}^M$, then derive $k^E_{AB},k^M_{AB}$\\
KDC performs access control as well, controlling with whom Alice can contact\\
Uses timestamp for freshness\\
If an attacker compromises the KDC, they have control over the master key and can therefore decrypt any message

\textbf{The GSM Handshake Protocol}\\
Mobile client, identified by $i$ (IMSI: International Mobile Subscriber Identifier)\\
Visited network (aka Base station), not fully trusted\\
Home network: trusted, shares key $k_i$ with client $i$\\
A5 function used to generate OTP is not a PRF (not secure)

\textbf{Attacks on GSM Handshake Protocol}\\
Visited network impersonation replay attack\\
Ciphersuite downgrade attack

\textbf{Visited-Network Impersonation Attack}\\
Cryptanalysis phase: attacker exposes $K_c$ based on the cyphertexts it collected in the eavesdropping phase\\
In the impersonate phase, the attacker will send the same $r$ and $s$ from before (replay attack), which will lead to the same $K_c$ it obtained in the cryptanalysis phase\\
GSM is using weak primitives that can be broken (OWF but not PRF A5/1, A5/2)\\
GSM doesn't check for freshness with $r$ or session key.

\textbf{GSM Ciphersuite Downgrade Attack}\\
Can be used to setup Visited-Network attack\\
A ciphersuite is the set of cryptographi cschemes used in protocol execution\\
Ciphersuite negotiation:\begin{itemize}
	\item Mobile sends a list of cipher-suites it supports
	\item Visited netwrk selected the strongest/most secure one that it also supports
	\item Goal of negotiation is to support interoperability b/w devices of different capabilities
\end{itemize}
A MitM attacker may trick these parties to use a weak suite (even if they can support a stronger one)\\
Works due to key reuse in GSM (same key is used across various encryption schemes)

\textbf{Improving Resiliency to Key Exposure}

\textbf{Forward Secrecy}\\
So far: session key $k_i^S \nRightarrow k_j^S$ (expose no other keys), and master key was fixed for all sessions.\\
Idea; Change the master key each session: $k_1^M,k_2^M,...$\\
Forward Secrecy (FS): \underline{master} key $k_i^M \nRightarrow k_j(j<i)$\\
I.e., $k_i^M$ (and $k_i^S$) don't expose keys, communication of \underline{previous} sessions $(j<i)$\\
Gives secrecy to the past (but not the future)

\textbf{Forward Secrecy 2PP Key Exchange}\\
This protocol generates a different master key for each session $i$ denoted as $k_i^M$

\end{document}
